\subsection{The amplitude and the global zero-Gram pair sum}
\label{sec:flag-global-amplitude}

Here we keep the actual row weights, normalized Poincare coefficients
and the column representation. The result below is the global
amplitude estimate and the summation of the fixed-label bound in the
original measures. We state the local-data hypothesis explicitly:
the transported local lowest vectors are those of the given
\(\mathrm{Choice}\ X\ 1\), with its \(\mathrm{WuBessel}\) and
\(\mathrm{GoodLoc}\) data. Thus the assertion is the amplitude
consumer of \(\mathrm{LocDatC}\), not a new construction of that
setting from abstract arithmetic parameters.

Write \(a_t=\|\operatorname{vec}(t)\|\) for the actual amplitude,
\(d_\pi=\dim\pi_u\) for the dimension of the original compact tensor
type, \(M=\operatorname{massRows}\), and \(\omega(t)=\operatorname{omI}(t)\)
for the infimum of the relative stabilizer indices over adapted normal
forms. Use the actual diagonal word \(u\), with the same norm \(T\)
for its two flags. Retain \(b,m,r,k,f,H=b^{C_Hm}\) and the
normalized minor labels of the preceding sections.

\paragraph{The positive rank condition.}
A relevant zero-Gram term outside the identity part has \(j\ge1\).
Indeed \(j=0\) in Lemma~\ref{lem:rational-flag-common-tail} means that
the full isotropic flags agree, with a common scalar \(c\) on all
graded lines and on the class of the last vector in \(X^\perp/X\).
The two last-vector classes are primitive in the same integral
quotient lattice (their coefficient ideals are both \(\mathcal O_E\)).
Consequently \(c\in\mathcal O_E^\times\). The last-vector norms give
\(c\bar c=1\), so Kronecker's theorem implies that \(c\) is a root of
unity. The term's upper-unitriangular admission at the two original
level primes gives \(z'-z\in X+g\Lambda_{\mathfrak p}\) for every
\(\mathfrak p\) above either prime. Pairing with \(z\) gives
\((c-1)T\in\mathfrak p\mathcal O_{E,\mathfrak p}\).
Since \(T\) is coprime to both level primes, \(c\equiv1\pmod{\mathfrak p}\).
A finite-order unit congruent to one at residue characteristic \(p\)
has \(p\)-power order: its prime-to-\(p\) part is also congruent to one,
and the factorization of \(x^n-1\) for prime-to-\(p\) order \(n\) shows
that this part is one. The two primes have different characteristics;
hence \(c=1\). The original two flags are now the same on every
successive graded line. Lemma~\ref{lem:integral-frame-movement}, applied
to the original admitted flags, supplies the original integral
unipotent flag equivalence. The term is in the identity part.
This proves \(\mathrm{JPosC}\), including the primitivity and
coprimality checks, without discarding any positive-rank term.

\begin{theorem}[Actual amplitude under the matched local-data hypothesis]
\label{thm:flag-global-amplitude}
Use the original diagonal term family and the local-data hypothesis
above, with the established all-length moment bounds and original
size conditions. For each \(C_s\ge0\) the constants can be fixed in
the original order, with \(C_H\) fixed before \(C_s\), so that every
relevant nonidentity zero-Gram term satisfies
\[
 a_t\le \mathrm{wt}(t)\sqrt{d_\pi}\,\omega(t)\,
           A(j,\delta),
 \tag{GA1}
\]
where \(\Delta=\Delta_j(\delta)>0\), \(S=S(j,\delta;H)\), and
\[
 A(j,\delta)=
  (C/k)^{j(k-1-j)f}
  \Delta^{-(k-1-j)/2}
  b^{C^2m^2fr^4} S^{-C_s k}.
 \tag{GA2}
\]
Here \(A\) is exactly \(\mathrm{k2a}\) of the source. In particular
this proves the original \(\mathrm{AmpAt}\) estimate, uniformly
under the stated local-data hypothesis.
\end{theorem}
\begin{proof}
We first verify the normalizing factors. Use the local group Haar
measures transported from the original product real group.
The finite column-space vector is supported in its one actual row
pair, and has norm
\[
 a_t=\mathrm{wt}(t)\,(\mathrm{scl}\,\mathrm{vol}(K_2))\,
                        \|\mathrm{OO}_t\|_{\rm HS}.
 \tag{GA3}
\]
The phase has modulus one, and the other factors here are precisely
the ones in the original expansion. Write \(V_N\) for (RH3). For
each place \(\tau\), write \(L_\tau=\int\|\Phi_\tau\|^2\) for the
positive finite local squared norm in the normalized local quotient
measure, \(d_\tau\) for its compact dimension, and \(c_\tau=L_\tau/d_\tau\).
The source's diagonal factor \(\mathrm{sclOf}\) is
\(\sqrt{c_\tau d_\tau/L_\tau}=1\); its operator \(J_\tau\) is the
unnormalized local orbital integral divided by \(c_\tau\).

The global unipotent and compact Haar measures differ from the
product normalized local subgroup measures by positive finite
constants \(c_N,c_H'\), respectively. The global Haar measure
on the real group is the product of the very same local group
measures. The uniqueness in the Weil formula shows that the global
quotient measure is \((c_Nc_H')^{-1}\) times the product local quotient
measure. Because the given seed is the tensor of the given local
lowest vectors and has squared norm one, Tonelli's theorem gives
\[
 \prod_\tau L_\tau=c_Nc_H',\qquad
 \prod_\tau d_\tau=d_\pi .
\]
The actual local Bessel functions are the unitary transports of
these lowest vectors, so this identity uses their masses, not
masses of independently renormalized functions. The given global
fundamental unipotent measure is \(c_N V_N\), and the global compact
mass is \(c_H'\). In the diagonal case the original definition is
\[
 c_0=\frac{\mathrm{vol}(K_2)c_NV_Nc_H'}{d_\pi},
 \qquad \mathrm{scl}=1/c_0 .
\]
In particular, with positive finite measures, we obtain the exact
identity
\[
       (\mathrm{scl}\,\mathrm{vol}(K_2))\,V_N
                 \prod_\tau c_\tau=1 .
 \tag{GA4}
\]
This supplies the source's \(\mathrm{ConstC}\).

Fix any admitted adapted normal form, and use its measure with
scalar one from Theorem~\ref{thm:flag-reference-haar}.
The original stabilizer of the orbital integrand is the conjugate
orientation of this group, with the measure transported by the
original return. The inverse-orientation measure calculation
in Theorem~\ref{thm:actual-local-degenerate-bound} preserves the
scalar one. Product Weil disintegration therefore identifies
the global stabilizer quotient with the product of these local
quotients, with no measure scalar. A relevant term has equal
characters on the common real stabilizer: the differential equality
on its rational Lie algebra extends to each real completion,
and exponentiation in the connected unipotent group gives equality
of the characters. Thus the rank-one integrand is invariant
under this group.

The original arithmetic stabilizer is a cocompact lattice, by
(RH1)--(RH2); take the indicated triangular fundamental set.
Weil disintegration of the original orbital integral first over
this compact stabilizer quotient and then over the real stabilizer
quotient gives its covolume times the quotient orbital integral.
These are legitimate integrals: the local theorem proves absolute
integrability, the product has finite norm integral by Tonelli, and
the stabilizer quotient has finite volume. Tensoring local matrix
coefficients and using Fubini gives the product of the local
unnormalized operators; if the source realizes the tensor by an
isometric embedding, this is an upper bound for its
Hilbert--Schmidt norm. Conjugation by the given unitary maps does
not change that norm. It follows from (GA3)--(GA4) that
\[
 a_t V_N
  \le \mathrm{wt}(t)\,
     \operatorname{covol}_{\lambda_R}\Gamma
                  \prod_\tau\|J_\tau\|_{\rm HS}.
 \tag{GA5}
\]
This is the original \(\mathrm{FactC}\), with the quotient
orientation as in its definition. It does not compare measures
on different arbitrary Haar normalizations.

Use \(\|J_\tau\|_{\rm HS}\le\sqrt{d_\tau}\|J_\tau\|_{\rm op}\)
and the actual local-product estimate
(Theorem~\ref{thm:actual-degenerate-product}). Its leading
minors are the original normalized minors: in the real frame
each pair of isotropic vectors is scaled by
\(\rho^{-(r-a)}\rho^{-(r-b)}\). Taking a leading determinant gives
\(\rho^{-2\sum_{a<l}(r-a)}|\tau(\delta_l)|\). Unitriangular
adaptation does not change the determinant.
Thus the source's product of local terms is bounded by
\[
 \sqrt{d_\pi}\,\kappa^f
 (16\pi e/k)^{j(k-1-j)f}\Delta^{-(k-1-j)/2}
 (8rH)^{80r^4f}S^{-C_sk}.
 \tag{GA6}
\]
The required moment is \(q=(\lceil C_s\rceil+2)kr\).
The established all-length theorem gives its original
simultaneous tail estimate after the original moment thresholds
are chosen. The size conditions \(k\ge40r^4\), \(k\ge5r\),
\(1\le j\le r\) give all hypotheses of the product bound.

Now (RH2)--(RH3) bound the prefactor from (GA5) by
\(\omega_N v^{k(r-j)f}L_v^2\). Multiplying by \(\kappa^f\) in
(GA6) gives
\[
      \omega_N\,2^{r^2f}v^{-kjf}L_v^2
        \le \omega_N\,2^{r^2f}L_v^2 .
\]
The source's \(\rho\le b\), \(v\ge1\), \(b\ge2\) and
\(m=2r+2\) imply
\[
 \log_b\!\left(2^{r^2f}L_v^2(8rH)^{80r^4f}\right)
 \le f\{r^2+12(r+1)^2+
                  80r^4(m+3+C_Hm)\}.
\]
Here \(8v\le b^3\) and \(8r\le b^{m+3}\).
The last expression is at most \(C^2m^2fr^4\) for
\(C\ge100(C_H+10)\); for instance
\(m+3+C_Hm\le(C_H+2)m\), and \(m\ge4\).
Increase \(C\) also to be at least \(16\pi e\) and
\(2(C_H+1)\). All these choices precede the variation
of the arithmetic data. Since \(\Delta\), \(j\) and \(S\)
are invariant under adaptation, we have proved (GA1)
with \(\omega_N\) for \emph{every} adapted form. Taking the
infimum gives exactly \(\omega(t)\), as required.
\end{proof}

\paragraph{Irrelevant terms.}
For the original orbital integral, failure of relevance means that
the differential of the ratio of its two unipotent characters on
the common stabilizer is nonzero. The resulting character is
nontrivial on the common real group and trivial on the original
arithmetic stabilizer (the finite-level phase compatibility).
Average the rank-one integrand on the arithmetic stabilizer
quotient. Right translation of the quotient multiplies the
integral of the character by its nontrivial value; invariance of
quotient Haar measure therefore forces zero. This argument applies
also outside the normal forms just used: a rational unipotent
group has cocompact integral lattice in its real points.
One may see this by a central series of rational vector-group
quotients, using full additive integral lattices and clearing
denominators in rational polynomial coordinates; the successive
quotients give compact fundamental sets. Absolute convergence
of the original period expansion, already proved, justifies
the averaging. This is the original irrelevance vanishing, not
an estimate of an absolute value after removing the character.

\begin{theorem}[The full weighted zero-Gram amplitude sum]
\label{thm:flag-global-amplitude-sum}
Under the same original local-data hypothesis as in
Theorem~\ref{thm:flag-global-amplitude}, choose \(C_s=C_p+4=20\).
The pair-amplitude sum over \emph{all} original diagonal
nonidentity zero-Gram pairs having equal row label is summable,
and, by increasing only the original standing thresholds,
\[
 \sum_{(t,t')\in\mathrm{allPairs}(\mathrm{degSet})}
                       a_t a_{t'}
   \le d_\pi M \exp\bigl(-\tfrac32 kf\log k\bigr).
 \tag{GA7}
\]
This estimate is before the source's finite Levi factor and
column-dimension comparison.
\end{theorem}
\begin{proof}
Let \(E(j,\delta)=\Delta^{k/2}b^{C_h^2m^2fr^4}S^{C_p k}\)
be the fixed-label bound (Theorem~\ref{thm:flag-global-pair-count}),
where \(C_p=16\) and \(C_h\) is fixed after \(C_H\).
For an arbitrary finite subset of the original pair set, terms
with zero amplitude contribute zero. The others are relevant
and have \(1\le j\le r\), as just proved. Partition them by
both labels. Multiplication of (GA1) for the two terms and
application of the fixed-label theorem gives
\[
 \sum_{\rm finite\ pairs}a_ta_{t'}
 \le d_\pi M
        \left(\sum_{j=1}^r\sum_{\delta}
                              A(j,\delta)E(j,\delta)\right)^2 ,
 \tag{GA8}
\]
where the inner sum can be over the finitely many labels of
this subset. No replacement of \(\mathrm{allPairs}\) by
a differently weighted pair system is made.

For any actual label,
\[
 AE=(C/k)^{j(k-1-j)f}
          \Delta^{(j+1)/2}
          b^{(C^2+C_h^2)m^2fr^4}S^{-4k}.
\]
Since \(\Delta\le b^{Cm^2f}S^2\) and \(S\ge1\),
the residual power of \(S\) is at most \(S^{j+1-4k}\),
which is at most \(S^{-3r}\) under the original rank
inequality. The actual distinct-label sum
(Theorem~\ref{thm:actual-normalized-minor-sum}) bounds
the sum of the latter by \(b^{C_c^2m^2fr}\),
where \(C_c=21+4C_H\). Hence
\[
 \sum_\delta AE
 \le (C/k)^{j(k-1-j)f} b^{B m^6 f},
 \qquad B=C^2+C_h^2+C_c^2+C+1 .
\]
All constants here are fixed before \(k,f,b\) vary.

Use the literal size relation \(C_l m^6\log b\le k\log k\).
We may demand \(C_l\ge16B\), and \(C_l\ge\max(32,C^{16})\).
The original \(m\ge C_l\), \(k\ge m^9\) give
\(k\ge32\) and \(\log C\le(\log k)/16\).
For \(1\le j\le r\le k-2\),
\(j(k-1-j)\ge k-2\ge15k/16\).
Consequently
\[
 \log\left(\sum_{j,\delta}AE\right)
 \le \log r-\frac{225}{256}kf\log k+\frac1{16}kf\log k
 \le -\frac34 kf\log k .
\]
The final inequality uses \(\log r\le\log k\le(kf\log k)/32\);
the numerical coefficient is in fact \(201/256\).
The estimate is valid for every finite selection of labels
and every finite subset of the original pairs. Taking the
supremum of nonnegative finite sums proves summability and (GA7).
The thresholds have been chosen in the original order.
\end{proof}

The final norm of the zero-Gram operator uses an additional
finite-dimensional comparison. In the exact original notation
\(\mathrm{Gam1}\) is the product of the finite Levi indices,
\(P_g=\prod_v(1+3/q_v)\), and \(d_u=\dim C_u\) is the full
column dimension (not \(d_\pi\)). Once the original Gram inequality
and the column comparison are supplied,
\[
 \|K''_{uu}\|_{\rm HS}^2
 \le P_g\,\mathrm{Gam1}\sum_{\rm pairs}a_ta_{t'},\qquad
 \mathrm{Gam1}\,M\,d_\pi\le d_u ,
 \tag{GA9}
\]
(GA7) immediately gives the literal source saving
\(d_u\exp(-kf\log k)\), because \(P_g\le2\) and
\(k\ge32\). We keep (GA9) explicit until its original finite
interfaces are verified. The local-data existence and generic
nonidentity returns, as well as the final retained arithmetic
theta Gram comparison, are not conclusions of (GA7).
